\subsection{Exact functoriality of the finite Hurwitz construction}

\begin{statement}{D23.F29}{Proved}{Full-GT covariance of a unique finite Hurwitz orbit}
\end{statement}
Let $G$ be a finite group, let $\mathcal C$ be one of its conjugacy
classes, and let
\[
 \mathcal N_{n,\mathcal C}=\mathcal C^n/G
\]
carry the left Hurwitz action
\[
 \sigma_i(g_1,\ldots,g_n)=
 (g_1,\ldots,g_i g_{i+1}g_i^{-1},g_i,\ldots,g_n).
\]
Let $H=\ker(\lambda:\widehat{GT}\to\widehat{\mathbb Z}^{\times})$.
There is a continuous action $V$ of $H$ on this entire finite set
such that, for the standard Drinfeld action $\Phi_h$ on
$\widehat B_n$,
\[
 V_h q(b)V_h^{-1}=q(\Phi_h(b))\qquad(b\in\widehat B_n).
\]
In particular any $B_n$-orbit whose cardinality occurs just once
among all the orbits in $\mathcal N_{n,\mathcal C}$ is preserved
by $V$. On such an orbit the same identity gives an actual
permutation conjugacy between $q\circ\Phi_h$ and $q$.
This applies to the degree-$1344$ orbit in D23.P17. It does not
assert any corresponding kernel stability for an arbitrary
linear representation of $B_n$.

\begin{proof}
The mixed braid group used here is
\[
 M_n=B_{n,1}:=\{b\in B_{n+1}:\pi(b)(n+1)=n+1\},
\]
where $\pi:B_{n+1}\to S_{n+1}$ is the permutation map.
Forgetting the distinguished last strand gives the split Artin
extension
\[
 1\longrightarrow N_n\longrightarrow M_n
 \mathop{\longrightarrow}^{p}B_n\longrightarrow1,
 \qquad N_n\simeq F_n.
\]
The section is the subgroup generated by
$\sigma_1,\ldots,\sigma_{n-1}$. The kernel is the fundamental
group of the plane minus the $n$ ordinary points, with meridians
\[
 t_i=\sigma_n\cdots\sigma_{i+1}\sigma_i^2
       \sigma_{i+1}^{-1}\cdots\sigma_n^{-1}
       \quad(1\leq i\leq n).
\]
Equivalently, putting
$\beta_i=\sigma_i^{-1}\cdots\sigma_{n-1}^{-1}\in B_n$
and $\beta_n=1$, the braid relations give
$t_i=\beta_i t_n\beta_i^{-1}$. Consequently $N_n$ is the
normal closure of $t_n=\sigma_n^2$ in $M_n$.

The same extension is exact after profinite completion. Here
is a justification of the only potentially delicate injection.
Finite-index characteristic subgroups $K$ of the finitely
generated free group $N_n$ are cofinal among its finite-index
subgroups. The $B_n$-action on $N_n/K$ has finite image, so
\[
 M_n\longrightarrow (N_n/K)\rtimes
             \operatorname{im}(B_n\to\operatorname{Aut}(N_n/K))
\]
is a finite quotient that embeds $N_n/K$. Thus the topology
induced on $N_n$ is its full profinite topology. Its completion
is therefore the closure $\widehat N_n\subset\widehat M_n$.
The quotient of $\widehat M_n$ by that closure is
$\widehat B_n$, by the universal property and the section.
Also $M_n$ has finite index in $B_{n+1}$, so its completion is
the open subgroup defined by the same permutation condition.

For $(\lambda,f)\in\widehat{GT}$, the standard full-GT formula is
\[
 \Phi(\sigma_1)=\sigma_1^\lambda,\qquad
 \Phi(\sigma_i)=f(\eta_i,\sigma_i^2)^{-1}
                   \sigma_i^\lambda f(\eta_i,\sigma_i^2),
 \quad
 \eta_i=\sigma_{i-1}\cdots\sigma_1^2\cdots\sigma_{i-1}.
\]
This is the profinite full-GT action, not merely its Galois
restriction. Since $\lambda$ is odd and both inputs of each
$f$ are pure, $\Phi$ induces the identity on $S_{n+1}$.
It preserves $\widehat M_n$. The formulas with $i<n$ use only
the standard section and restrict there to the standard
$\Phi$ on $\widehat B_n$.

Put $u=f(\eta_n,t_n)$. Its image under $p$ is
$f(p(\eta_n),1)=1$, since $f$ is in the commutator subgroup
of the free profinite group. Thus $u\in\widehat N_n$, and
\[
 \Phi(t_n)=u^{-1}t_n^\lambda u\in\widehat N_n.
\]
Normal generation shows that $\Phi(\widehat N_n)$ is contained
in $\widehat N_n$. Applying the same argument to the inverse
GT element gives equality. Write
$\alpha=\Phi|_{\widehat N_n}$. Preservation of the section
gives the exact compatibility
\[
 \alpha\theta_b\alpha^{-1}=\theta_{\Phi(b)},
 \qquad \theta_b(t)=btb^{-1}.
\]

The Artin action of a braid $b\in B_n$ takes $t_i$ to an
$N_n$-conjugate of $t_{\pi(b)(i)}$. This follows directly from
the two-generator formulas
\[
 \theta_{\sigma_i}^{-1}(t_i)=t_i t_{i+1}t_i^{-1},\qquad
 \theta_{\sigma_i}^{-1}(t_{i+1})=t_i,
\]
with all other meridians fixed. The assertion extends to
$\widehat B_n$: the conjugacy class of an element in the
compact group $\widehat N_n$ is closed, and discrete braids
are dense in each fixed permutation coset. Since
$\Phi(\beta_i)$ has the same permutation as $\beta_i$, the
formula for $\Phi(t_n)$ and $t_i=\beta_i t_n\beta_i^{-1}$
show that $\alpha(t_i)$ is an $\widehat N_n$-conjugate of
$t_i^\lambda$. In particular for $h\in H$ each peripheral
conjugacy class is preserved.

Identify $\mathcal C^n$ with continuous homomorphisms
$\varphi:\widehat N_n\to G$ taking every $t_i$ into
$\mathcal C$. Quotienting by simultaneous $G$-conjugation
gives $\mathcal N_{n,\mathcal C}$. With the meridian convention
just fixed, the literal Hurwitz action in the statement is
\[
 q(b)[\varphi]=[\varphi\circ\theta_b^{-1}].
\]
Define
\[
 V_h[\varphi]=[\varphi\circ\alpha_h^{-1}].
\]
The peripheral result makes this well defined on the entire
specified tuple set. It is a bijection, respects simultaneous
conjugation, and is a continuous action of $H$: the set of
homomorphisms is finite and evaluation of the images of the
finitely many generators is continuous. Exact compatibility
gives
\[
 \begin{split}
 V_h q(b)[\varphi]
   &=[\varphi\circ\theta_b^{-1}\circ\alpha_h^{-1}]\\
   &=[\varphi\circ\alpha_h^{-1}
                  \circ\theta_{\Phi_h(b)}^{-1}]
     =q(\Phi_h(b))V_h[\varphi].
 \end{split}
\]
Since $\Phi_h$ is an automorphism of $\widehat B_n$, $V_h$
permutes its orbits and preserves their sizes. The uniqueness
criterion follows. D23.P17 checks uniqueness over all of
$\mathcal C^4/G$, not over a subset selected by product class.
Finally the full twist acts on a tuple by simultaneous
conjugation by its product, so the action on this quotient
already factors through $B_n/Z(B_n)$, as asserted there.
\end{proof}

\paragraph{Primary source and scope.}
Harbater--Schneps, \emph{Fundamental groups of moduli and the
Grothendieck--Teichm\"uller group}, author PDF
\url{https://webusers.imj-prg.fr/~leila.schneps/HS.pdf},
Section~0.1, equation~(1), printed page~3 (PDF page index~2),
explicitly gives the profinite full-GT braid action above.
Section~3, Lemma~1(a), printed page~30 (PDF index~29), also
records the pure peripheral conjugacy property on the sphere
mapping-class quotient. The mixed Artin extension and its
profinite exactness, as needed here, have been justified in
the proof. Nakamura, RIMS1691, Section~5.4, equation~(5.4.2),
printed page~28, supplies the arithmetic tangential marking,
but that passage alone is Galois-scoped and is not used to
upgrade a Galois action to all of $\widehat{GT}$.

\paragraph{Immediate finite constraints.}
For D23.P17 write $A=q(\sigma_1)$, $B=q(\sigma_2)$,
$C=q(\sigma_3)$, $F=q(f(\sigma_1^2,\sigma_2^2))$,
and $R=q(f(\sigma_3^2,(\sigma_1\sigma_2)^3))$.
Then every actual $h=(1,f)\in\widehat{GT}$ admits the same
permutation $V_h$ satisfying
\[
 V_h A V_h^{-1}=A,\quad
 V_h B V_h^{-1}=F^{-1}BF,\quad
 V_h C V_h^{-1}=RCR^{-1}.
\]
For the last identity, the central factor $\sigma_1^2$
in $(\sigma_1\sigma_2)^3=\sigma_1^2\eta_3$ commutes with
both arguments of $f$, so the $\eta_3$ input can be replaced
by that full twist; the two-cycle gives the displayed
orientation. These are necessary finite conditions for actual
full GT. No arbitrary tuple satisfying them is thereby
promoted to a gentle or full-GT element.
